\documentclass[11pt,a4paper]{article}
\usepackage[utf8]{inputenc}
\usepackage[margin=1in]{geometry}
\usepackage{amsmath,amssymb}
\usepackage{graphicx}
\usepackage{booktabs}
\usepackage{hyperref}
\usepackage{listings}
\usepackage{xcolor}
\usepackage{microtype}
\usepackage{caption}

\hypersetup{
    colorlinks=true,
    linkcolor=blue,
    filecolor=magenta,      
    urlcolor=cyan,
    citecolor=blue
}

\lstset{
    basicstyle=\ttfamily\small,
    breaklines=true,
    frame=single,
    backgroundcolor=\color{gray!10},
    keywordstyle=\color{blue}\bfseries,
    commentstyle=\color{green!50!black}\itshape,
    stringstyle=\color{red!60!black}
}

\title{\textbf{Real-Time Chat \& Streaming Analytics Platform}\\
\large Architecture, Implementation, and Evaluation with Apache Kafka and Spark Structured Streaming}
\author{Big Data Engineering Team}
\date{\today}

\begin{document}

\maketitle

\begin{abstract}
Modern messaging systems require not only low-latency interactive communication, but also real-time visibility into traffic volume, user engagement, and conversation dynamics. Traditional operational databases (OLTP) degrade significantly when queried for live rolling metrics. In this project, we designed and implemented a dual-path Big Data platform that separates low-latency operational chat flows from high-throughput streaming analytics. By integrating FastAPI and WebSocket with Apache Kafka (KRaft mode) as an ingestion buffer, Apache Spark Structured Streaming for tumbling window aggregations with late-data watermarking, and PostgreSQL for decoupled analytical query serving, our system achieves sub-second chat latency while delivering updated aggregated metrics every few seconds. We demonstrate the stability and throughput of the platform under simulated load using a dedicated traffic generation pipeline.
\end{abstract}

\section{Introduction}
Real-time chat platforms like Slack, Discord, and Telegram experience dynamic, bursty traffic patterns. While traditional database architectures can satisfy point-to-point operational messaging, computing rolling metrics (such as messages per minute, distinct active users, and trending conversations) directly against transactional OLTP tables causes severe lock contention, IO bottlenecks, and degraded user experience.

To solve this, Big Data architectures decouple the analytical stream from the transactional write path. This paper presents the architecture, engineering design, and experimental evaluation of the \textbf{Real-Time Chat \& Streaming Analytics Platform}.

\section{System Architecture}
The system architecture adheres strictly to the single-responsibility principle across components:

\begin{enumerate}
    \item \textbf{Frontend Client}: A modern React 19 single-page application built with Vite and TailwindCSS, offering responsive chat messaging alongside a live streaming analytics dashboard powered by Recharts.
    \item \textbf{API \& Ingestion Gateway}: FastAPI backend handling WebSocket connections, authentication, and HTTP REST APIs. Upon successfully committing a message to PostgreSQL (OLTP), it publishes an event asynchronously to Apache Kafka.
    \item \textbf{Message Ingestion Layer}: Apache Kafka running in KRaft mode (Kafka Raft metadata mode without Apache ZooKeeper), configured with topic \texttt{chat-events} partitioned by \texttt{conversation\_id}.
    \item \textbf{Stream Processing Layer}: Apache Spark Structured Streaming computing tumbling window metrics with event-time semantics, watermarking, and idempotent database UPSERT sinks.
    \item \textbf{Decoupled Storage Layer}:
    \begin{itemize}
        \item \textbf{PostgreSQL (OLTP)}: Houses business operational data (users, messages, friendships, attachments).
        \item \textbf{PostgreSQL (Analytics)}: Houses aggregated tables (\texttt{stream\_messages\_per\_minute}, \texttt{stream\_active\_users\_per\_minute}, \texttt{stream\_messages\_per\_conversation}).
        \item \textbf{Redis 7}: Serves presence detection, typing indicators, token caching, and WebSocket horizontal pub/sub broadcasting.
    \end{itemize}
\end{enumerate}

\section{Event Streaming \& Partitioning with Apache Kafka}
Kafka acts as an asynchronous shock absorber between the real-time chat application and the analytical engine.

\subsection{Event Schema}
Each chat event conforms to the following standardized JSON format:
\begin{lstlisting}[language=json]
{
  "event_type": "message_sent",
  "message_id": "8f1a23c4-...",
  "user_id": "3b4c5d6e-...",
  "conversation_id": "7a8b9c0d-...",
  "message_type": "text",
  "timestamp": "2026-09-28T13:20:00Z"
}
\end{lstlisting}

\subsection{Partitioning Strategy}
Events are partitioned on Kafka using \texttt{conversation\_id} as the record key. This guarantees total order of events for any individual conversation across partitions while enabling horizontal scaling across consumer cores.

\subsection{Fault Isolation \& Fire-and-Forget}
Kafka publishing is performed \emph{after} the operational PostgreSQL commit is completed. The producer publishes events in a non-blocking, fire-and-forget fashion. Should Kafka experience temporary unavailability or backpressure, error logs are emitted, but end-user chat operations remain uninterrupted.

\section{Stream Analytics with Apache Spark Structured Streaming}
The analytical engine is built with PySpark Structured Streaming.

\subsection{Event-Time Windowing and Watermarking}
In distributed chat environments, network variability can cause late arrival of messages. The pipeline extracts the event timestamp and enforces a \textbf{30-second watermark}:
\begin{lstlisting}[language=Python]
watermarked_stream = parsed_stream \
    .withWatermark("event_time", "30 seconds")
\end{lstlisting}
Events arriving within 30 seconds of the latest observed timestamp are incorporated into their respective \textbf{1-minute tumbling windows}. Events trailing beyond 30 seconds are safely dropped to conserve executor state memory.

\subsection{Analytical Metrics}
Three concurrent streaming pipelines execute across micro-batches:
\begin{enumerate}
    \item \textbf{Messages Per Minute}: Aggregates total message counts per 1-minute window:
    \[
    M(t) = \sum_{e \in W_t} 1
    \]
    \item \textbf{Active Users Per Minute}: Calculates distinct contributors per window via \texttt{countDistinct(user\_id)}:
    \[
    U(t) = |\{e.\text{user\_id} : e \in W_t\}|
    \]
    \item \textbf{Messages Per Conversation}: Groups message counts by $(W_t, \text{conversation\_id})$.
\end{enumerate}

\subsection{Idempotent Micro-Batch Sinks (foreachBatch UPSERT)}
To guard against duplicates resulting from at-least-once delivery semantics and executor retries, Spark writes records to PostgreSQL using \texttt{foreachBatch} combined with atomic SQL \texttt{ON CONFLICT DO UPDATE} (UPSERT) clauses:
\begin{lstlisting}[language=SQL]
INSERT INTO stream_messages_per_minute (window_start, window_end, message_count)
VALUES (%s, %s, %s)
ON CONFLICT (window_start, window_end)
DO UPDATE SET message_count = EXCLUDED.message_count;
\end{lstlisting}

\section{Analytics Serving \& Frontend Dashboard}
The analytical tables are exposed via lightweight, indexed FastAPI endpoints:
\begin{itemize}
    \item \texttt{GET /api/v1/analytics/messages-per-minute}
    \item \texttt{GET /api/v1/analytics/active-users}
    \item \texttt{GET /api/v1/analytics/conversations}
\end{itemize}
The frontend queries these endpoints periodically (configurable 5-second polling) and visualizes the streaming metrics via an interactive Recharts dashboard featuring gradient area charts, bar charts, and micro-batch audit tables.

\section{Experimental Evaluation}
A traffic generator script (\texttt{producer\_test.py}) was utilized to benchmark the streaming pipeline under varying event injection rates ranging from $10$ to $50$ messages per second across multiple concurrent mock users and conversations.

\begin{table}[h]
\centering
\caption{System Performance Under Simulated Traffic}
\begin{tabular}{@{}lccc@{}}
\toprule
Metric & 10 msgs/sec & 25 msgs/sec & 50 msgs/sec \\
\midrule
FastAPI Kafka Publish Latency & $<2.5$ ms & $<4.1$ ms & $<6.8$ ms \\
Spark Micro-batch Trigger Duration & 5.0 s & 5.0 s & 5.0 s \\
End-to-End Analytics Latency & $5.8 \pm 1.2$ s & $6.2 \pm 1.4$ s & $6.9 \pm 1.7$ s \\
PostgreSQL Upsert Batch Time & $<35$ ms & $<55$ ms & $<90$ ms \\
\bottomrule
\end{tabular}
\end{table}

The empirical findings confirm that the system handles bursty loads efficiently without causing backpressure or degradation on the transactional chat pathway.

\section{Conclusion}
The Real-Time Chat \& Streaming Analytics Platform demonstrates how Big Data technologies can be integrated cleanly into modern web applications. By utilizing Apache Kafka as an asynchronous buffer and Apache Spark Structured Streaming for windowed computation, real-time analytics are generated accurately with bounded latency while maintaining total isolation from transactional OLTP operations.

\end{document}
